\documentclass[11pt,a4paper]{article}

\usepackage[a4paper,margin=2.2cm]{geometry}
\usepackage[T1]{fontenc}
\usepackage[utf8]{inputenc}
\usepackage{lmodern}
\usepackage{microtype}
\usepackage{enumitem}
\usepackage{array}
\usepackage{booktabs}
\usepackage{tabularx}
\usepackage{titlesec}
\usepackage[table]{xcolor}
\usepackage{graphicx}
\usepackage{fancyhdr}
\usepackage{hyperref}
\usepackage{multirow}

\definecolor{kmutnbBlue}{HTML}{1F3A8A}
\definecolor{accent}{HTML}{0F6FC5}
\definecolor{lightGray}{HTML}{F2F4F7}

\hypersetup{
    colorlinks=true,
    linkcolor=accent,
    urlcolor=accent,
    citecolor=accent
}

\titleformat{\section}
  {\normalfont\large\bfseries\color{kmutnbBlue}}
  {\thesection}{0.6em}{}
\titlespacing*{\section}{0pt}{1.2\baselineskip}{0.4\baselineskip}

\setlist[itemize]{leftmargin=1.4em,itemsep=2pt,topsep=2pt}
\setlist[enumerate]{leftmargin=1.6em,itemsep=2pt,topsep=2pt}

\pagestyle{fancy}
\fancyhf{}
\fancyhead[L]{\small\textit{010153103 Algorithms and Data Structures}}
\fancyhead[R]{\small Semester 1/2026}
\fancyfoot[C]{\small\thepage}
\renewcommand{\headrulewidth}{0.3pt}

\newcommand{\fieldlabel}[1]{\textbf{#1:}}

\begin{document}

%% ─── Title block ────────────────────────────────────────────────────────────
\begin{center}
  {\Large\bfseries\color{kmutnbBlue}
    Electrical and Computer Engineering Course Syllabus}\\[4pt]
  {\large\bfseries 010153103 \quad Algorithms and Data Structures}\\[2pt]
  {\normalsize Semester 1 / Academic Year 2026}\\[2pt]
  {\small Department of Electrical and Computer Engineering (ECE)}\\
  {\small Faculty of Engineering,
    King Mongkut's University of Technology North Bangkok}\\[4pt]
  {\small Instructor: Asst.\ Prof.\ Dr.\ Ruslee Sutthaweekul
    \quad\textbar\quad
    \href{mailto:ruslee.s@eng.kmutnb.ac.th}{ruslee.s@eng.kmutnb.ac.th}}
\end{center}

\noindent\rule{\linewidth}{0.4pt}\medskip

%% ─── 1. Course Information ──────────────────────────────────────────────────
\section*{1.\quad Course Information}
\begin{tabular}{@{}p{4.4cm}p{10.8cm}@{}}
  \fieldlabel{Course Number}   & 010153103 \\[2pt]
  \fieldlabel{Course Name}     & Algorithms and Data Structures \\[2pt]
  \fieldlabel{Credits}         & 3 (3--0--6) \\[2pt]
  \fieldlabel{Contact Hours}   & 3 hours per week \\[2pt]
  \fieldlabel{Course Type}     & Required \\[2pt]
  \fieldlabel{Prerequisites}   & 010153002 Computer Programming \\[2pt]
  \fieldlabel{Textbook}        & Goodrich, M.~T., Tamassia, R., \& Goldwasser, M.~H.\
                                  \textit{Data Structures and Algorithms in Java},
                                  6th Edition. John Wiley \& Sons, 2014. \\
\end{tabular}

%% ─── 2. Course Description ──────────────────────────────────────────────────
\section*{2.\quad Course Description}
This course requires students to understand basic data structures for computer
programming development, including \emph{lists}, \emph{stacks}, \emph{queues},
and \emph{trees}; core algorithmic techniques for \emph{sorting},
\emph{searching}, and \emph{hashing}; and the use of \emph{recursion}.
Students will learn to select and apply appropriate data structures and
algorithms to solve specific programming problems.

%% ─── 3. Course Learning Outcomes ────────────────────────────────────────────
\section*{3.\quad Course Learning Outcomes (CLOs)}
Upon successful completion of this course, students will be able to:
\begin{enumerate}
  \item Demonstrate a good understanding of linear data structures, including
        arrays, linked lists, stacks, and queues.
  \item Implement abstract data types effectively.
  \item Understand non-linear data structures, including trees and graphs.
  \item Select and use appropriate data structures to solve specific problems.
  \item Understand and implement sorting, searching, and hashing algorithms.
  \item Choose and apply appropriate algorithms to solve specific problems.
\end{enumerate}

%% ─── 4. ABET Student Outcomes ───────────────────────────────────────────────
\section*{4.\quad ABET Student Outcomes Addressed}
\renewcommand{\arraystretch}{1.35}
\begin{tabularx}{\linewidth}{X >{\centering\arraybackslash}p{3.2cm}}
  \toprule
  \rowcolor{kmutnbBlue!90}
  \textcolor{white}{\textbf{ABET Student Outcome (SO)}} &
  \textcolor{white}{\textbf{Course Learning Outcome (CLO)}} \\
  \midrule
  1.\ An ability to identify, formulate, and solve complex engineering problems
      by applying principles of engineering, science, and mathematics.
  & 1--6 \\
  \rowcolor{lightGray}
  2.\ An ability to apply engineering design to produce solutions that meet
      specified needs with consideration of public health, safety, and welfare,
      as well as global, cultural, social, environmental, and economic factors.
  & 4, 6 \\
  3.\ An ability to communicate effectively with a range of audiences.
  & --- \\
  \rowcolor{lightGray}
  4.\ An ability to recognize ethical and professional responsibilities in
      engineering situations and make informed judgments, which must consider
      the impact of engineering solutions in global, economic, environmental,
      and societal contexts.
  & --- \\
  5.\ An ability to function effectively on a team whose members together
      provide leadership, create a collaborative and inclusive environment,
      establish goals, plan tasks, and meet objectives.
  & --- \\
  \rowcolor{lightGray}
  6.\ An ability to develop and conduct appropriate experimentation, analyze
      and interpret data, and use engineering judgment to draw conclusions.
  & 2, 4, 6 \\
  7.\ An ability to acquire and apply new knowledge as needed, using
      appropriate learning strategies.
  & 4, 6 \\
  \bottomrule
\end{tabularx}

%% ─── 5. CLO–Assessment Mapping ──────────────────────────────────────────────
\section*{5.\quad CLO Assessment Mapping}
\renewcommand{\arraystretch}{1.3}
\begin{tabularx}{\linewidth}{X >{\centering\arraybackslash}p{2.0cm} >{\centering\arraybackslash}p{3.8cm}}
  \toprule
  \rowcolor{kmutnbBlue!90}
  \textcolor{white}{\textbf{Course Learning Outcome (CLO)}} &
  \textcolor{white}{\textbf{SO}} &
  \textcolor{white}{\textbf{Assessment Tools}} \\
  \midrule
  1.\ Demonstrate a good understanding of linear data structures,
      including arrays, linked lists, stacks, and queues.
  & 1 & Exams \\
  \rowcolor{lightGray}
  2.\ Implement abstract data types effectively.
  & 1, 6 & Exams, Assignment \\
  3.\ Understand non-linear data structures, including trees and graphs.
  & 1 & Exams \\
  \rowcolor{lightGray}
  4.\ Select and use appropriate data structures to solve specific problems.
  & 1, 2, 6, 7 & Exams, Assignment \\
  5.\ Understand and implement sorting, searching, and hashing algorithms.
  & 1 & Exams \\
  \rowcolor{lightGray}
  6.\ Choose and apply appropriate algorithms to solve specific problems.
  & 1, 2, 6, 7 & Exams, Assignment \\
  \bottomrule
\end{tabularx}

%% ─── 6. Weekly Schedule ─────────────────────────────────────────────────────
\section*{6.\quad Weekly Schedule}
\renewcommand{\arraystretch}{1.25}
\rowcolors{2}{lightGray}{white}
\begin{tabularx}{\linewidth}{>{\centering\arraybackslash}p{1.3cm} X >{\centering\arraybackslash}p{2.6cm}}
  \toprule
  \rowcolor{kmutnbBlue!90}
  \textcolor{white}{\textbf{Week}} &
  \textcolor{white}{\textbf{Topic}} &
  \textcolor{white}{\textbf{Reading}} \\
  \midrule
  1  & Introduction to Algorithms \& Data Structures; Programming Review   & Ch.\ 1   \\
  2  & Java \& Object-Oriented Design                                       & Ch.\ 2   \\
  3  & Algorithm Analysis (Big-O, Big-$\Omega$, Big-$\Theta$)              & Ch.\ 4   \\
  4  & Array Lists                                                          & Ch.\ 3   \\
  5  & Singly Linked Lists                                                  & Ch.\ 3   \\
  6  & Doubly Linked Lists                                                  & Ch.\ 3   \\
  7  & Stacks \& Queues                                                     & Ch.\ 6   \\
  \rowcolor{accent!15}
  \multicolumn{1}{c}{---} & \textbf{Midterm Examination}                   & ---      \\
  8  & Recursion                                                            & Ch.\ 5   \\
  9  & Trees \& Binary Trees                                                & Ch.\ 8   \\
  10 & Graphs                                                               & Ch.\ 14  \\
  11 & Sorting Algorithms (Merge Sort, Quick Sort, Heap Sort)               & Ch.\ 12  \\
  12 & Searching Algorithms (Linear, Binary)                                & Ch.\ 12  \\
  13 & Hashing (Hash Tables, Collision Resolution)                          & Ch.\ 10  \\
  14 & Review \& Problem-Solving Practice                                   & ---      \\
  \rowcolor{accent!15}
  \multicolumn{1}{c}{---} & \textbf{Final Examination}                     & ---      \\
  \bottomrule
\end{tabularx}
\rowcolors{1}{}{}

\medskip
\noindent\textit{Note:} The schedule above is subject to minor adjustments
to accommodate the academic calendar and class progress.

%% ─── 7. Assessment & Grading ────────────────────────────────────────────────
\section*{7.\quad Assessment and Grading}

\noindent
\begin{minipage}[t]{0.48\linewidth}
  \textbf{Assessment Breakdown}
  \begin{itemize}
    \item Midterm Examination \dotfill 30\%
    \item Final Examination   \dotfill 50\%
    \item Assignments         \dotfill 20\%
  \end{itemize}
\end{minipage}\hfill
\begin{minipage}[t]{0.48\linewidth}
  \textbf{Grading Scale}\par\smallskip
  \begin{tabular}{@{}l@{\hspace{1em}}l@{\hspace{2.5em}}l@{\hspace{1em}}l@{}}
    81--100 & A  & 51--60 & C  \\
    76--80  & B+ & 46--50 & D+ \\
    71--75  & B  & 41--45 & D  \\
    61--70  & C+ & 0--40  & F  \\
  \end{tabular}
\end{minipage}

%% ─── 8. Course Materials ────────────────────────────────────────────────────
\section*{8.\quad Course Materials}
\begin{itemize}
  \item Lecture slides: \href{https://ruslylove.github.io/datastruct/}{ruslylove.github.io/datastruct}
  \item PDF export of slides: \href{https://ruslylove.github.io/datastruct/slides.pdf}{slides.pdf}
  \item Source repository: \href{https://github.com/ruslylove/datastruct}{github.com/ruslylove/datastruct}
\end{itemize}

%% ─── Footer ─────────────────────────────────────────────────────────────────
\vfill
\noindent\rule{\linewidth}{0.4pt}
\begin{center}
  \small\textit{This syllabus is subject to minor adjustments at the
  instructor's discretion to accommodate the academic calendar and class
  progress. Students are responsible for staying informed of any changes
  announced in class or via the course website.}
\end{center}

\end{document}
